\chapter{First Contact with CP-SAT}\label{chap:5}

\begin{quote}
\emph{``Make it work. Make it right. Make it fast.''} --- Kent Beck
\end{quote}

The solver this manuscript commits to is \textbf{CP-SAT}, the
constraint-programming solver that ships inside Google's
\href{https://developers.google.com/optimization}{OR-Tools}. It is open
source, distributed under the Apache 2.0 license, and free for any use.
Installation is a single \lstinline!pip! command, runs on
Linux, macOS, and Windows, and requires no license server or external
dependencies.

The primary reason for this choice is that CP-SAT occupies a useful
middle ground for learning. Its constraint vocabulary is rich --- global
constraints like \lstinline!circuit!,
\lstinline!no_overlap!, and
\lstinline!all_different! let you express combinatorial
structures directly, without the forced linearization a pure MIP solver
would require. But CP-SAT's Python interface is still a
\emph{solver-specific API}, not an algebraic modeling language that
targets multiple backends behind an abstraction layer. You are building
a model that a specific solver will run, and you see how each modeling
decision maps to what the solver actually works with. That directness is
valuable when learning: it keeps the connection between the model and
the solver visible.

CP-SAT has won the \href{https://www.minizinc.org/challenge/}{MiniZinc
Challenge}, the standard annual benchmark in the
constraint-programming community, every year since 2018. For
combinatorial problems with discrete decisions, it is one of the
strongest tools available.

One property worth stating up front: CP-SAT is \textbf{integer-only}.
Booleans, finite enumerations, integer-valued amounts, and
integer-modeled time variables are all native. Continuous quantities are
not supported and must be either integerized (multiply real prices in
dollars by 100 and work in cents) or pushed to a different solver
(linear programming, mixed-integer programming with continuous
variables). For genuinely combinatorial problems the integer restriction
is rarely felt; for problems with a continuous core, CP-SAT is the wrong
tool.

\begin{platypusbook}[The CP-SAT Primer]
This chapter and the two reference chapters that follow cover just enough of
CP-SAT to model with it. The companion
\href{https://github.com/d-krupke/cpsat-primer}{CP-SAT Primer} is the
comprehensive treatment: it walks through the full feature set, the solver
parameters, and the performance engineering this manuscript sets aside. Reach
for it whenever you want more depth than a modeling text needs to give.
\end{platypusbook}

\section{What Is Inside the Box}\label{sec:5.1}

A reader who is new to constraint solvers might reasonably expect CP-SAT
to be one algorithm. It is not. CP-SAT is a \emph{portfolio}: several
complementary techniques wired together, run in parallel, and sharing
learned information with each other.

A short tour of what is inside:

\begin{itemize}
\item
  A \textbf{SAT} core handles Boolean reasoning, using the
  conflict-driven clause-learning algorithm that powers modern SAT
  solvers. This is what makes CP-SAT extraordinarily good at logical
  constraints.
\item
  A \textbf{CP} layer adds the global constraints,
  \lstinline!all_different!,
  \lstinline!circuit!,
  \lstinline!cumulative!, and others, with dedicated
  propagators that reason about whole structures at once rather than
  constraint by constraint.
\item
  A \textbf{linear-programming relaxation} runs alongside to provide
  tight dual bounds. This is what allows CP-SAT to make serious inroads
  into territory traditionally dominated by MIP solvers.
\item
  \textbf{Large Neighborhood Search} (LNS) takes a current solution,
  freezes most of it, and re-optimizes the rest, a way of escaping
  local minima that hand-rolled procedures rarely match.
\item
  An aggressive \textbf{presolve} simplifies the model before search
  begins, often removing variables and constraints the modeler did not
  realize were redundant.
\item
  Several \textbf{workers} run different combinations of these
  techniques in parallel and share learned clauses, so a discovery in
  one worker can prune the search of another.
\end{itemize}

The crucial implication for the modeler: \emph{you typically do not pick
the strategy explicitly}. You hand the solver a model and call
\lstinline!solve!; the portfolio races itself for the best
answer. Advanced usage does involve steering (branching strategies,
LNS parameters, worker counts, time limits, and so on), but most of
the leverage on real problems comes from the model itself rather than
from solver tuning; for the cases where tuning matters, later chapters
discuss it. For now, ``build a model, call
\lstinline!solve!, get an answer'' is a complete mental
model.

\section{Where CP-SAT Fits, and Where It Does Not}\label{sec:5.2}

CP-SAT is opinionated. It excels in a particular shape of problem and is
the wrong tool in others. Knowing the difference is part of using it
well.

\paragraph{Where CP-SAT fits.} The sweet spot is problems whose decisions
are dominantly discrete and whose constraints are richly logical:

\begin{itemize}
\item
  \textbf{Scheduling and rostering}: assigning tasks to machines,
  workers to shifts, courses to rooms.
\item
  \textbf{Routing and tour problems}: vehicle routing on small to
  medium graphs, sequencing.
\item
  \textbf{Packing, assignment, and matching}: bin packing,
  generalized assignment, stable matching with side constraints.
\item
  \textbf{Discrete planning} with rich logical rules, especially when
  the rules are messy in ways that pure linear formulations do not
  handle gracefully.
\item
  \textbf{Feasibility checking}: \emph{is there any valid
  configuration at all?} CP-SAT is exceptionally good at this; problems
  with hundreds of logical rules that look intractable are often solved
  or proven infeasible in seconds.
\end{itemize}

\paragraph{Where to reach for something else.} CP-SAT is not the right tool
when:

\begin{itemize}
\item
  Continuous variables dominate the problem. Reach for a
  linear-programming or nonlinear-programming solver. (Cents-and-dollars
  integerization works for shallow continuous components; deep
  continuous structure does not survive it.)
\item
  The problem is purely linear, the LP relaxation is strong, and the
  scale is large. A dedicated mixed-integer programmer such as Gurobi,
  CPLEX, or HiGHS often wins on the standard benchmarks. CP-SAT has an
  LP layer of its own, but it is not optimized for the regimes the
  dedicated MIP solvers are.
\item
  The problem is the symmetric Traveling Salesman Problem at very large
  scale (tens of thousands of cities). Specialist solvers like Concorde
  and LKH remain ahead.
\item
  The objective is differentiable and you want gradients. JAX, PyTorch,
  and the gradient-based optimization stack are an entirely different
  shape of tool.
\end{itemize}

CP-SAT thrives on combinatorial structure with rich logical constraints;
mixed-integer programming thrives on polyhedral strength and tight
linear relaxations. Many real problems contain both flavors, in which
case the right tool comes down to which flavor dominates.

The rest of this manuscript works inside CP-SAT's sweet spot.

\section{Setup}\label{sec:5.3}

\begin{lstlisting}[language=bash]
pip install ortools
\end{lstlisting}

That is the entire installation. The package ships precompiled wheels
for the major platforms (Linux, macOS, Windows, on x86-64 and ARM)
so there is no C++ build step on a fresh machine. Python 3.9 and newer
are supported.

Two imports are used in essentially every CP-SAT program:

\begin{lstlisting}[language=Python]
from ortools.sat.python import cp_model
\end{lstlisting}

\lstinline!cp_model.CpModel! is the model \emph{builder},
the object you populate with variables and constraints.
\lstinline!cp_model.CpSolver! is the \emph{solver
wrapper}, the object that takes a built model, runs the search, and
exposes the answer. Every example in the rest of the manuscript starts
with this import.

A note on hardware: a modern laptop is sufficient to get started. CP-SAT
does not use a GPU. The two resources that matter are cores (CP-SAT runs
multiple workers in parallel) and memory (presolve, LP relaxation, and
learned clauses all need RAM). Four cores and 16 GB of RAM is a
reasonable starting point for learning, but serious projects usually
benefit from more: both more workers and more memory headroom.
Smaller machines are still usable for modest instances.

\section{A Complete Example: Facility Location}\label{sec:5.4}

We walk through the \textbf{uncapacitated facility location problem}
(FLP) end to end, with the full five-step recipe from \Cref{chap:3} visible
in the code. The FLP is a classical assignment problem with an
objective: pick which warehouses to open and which warehouse serves each
customer.

The problem statement:

\begin{quote}
A delivery company is planning to open new warehouses to serve its
customers. There is a set \(F\) of potential warehouse locations and a
set \(C\) of customers who need to receive their orders. Opening a
warehouse at location \(i \in F\) comes with a fixed cost \(f_i\).
Shipping goods from warehouse \(i\) to customer \(j \in C\) costs
\(c_{i,j}\). Every customer must be served by exactly one warehouse, and
a warehouse can only ship goods if it is actually opened. The company
needs to decide which warehouse locations to open and which warehouse
will serve each customer so that the total cost, opening costs plus
shipping costs, is as small as possible.
\end{quote}

We apply the five-step recipe from \Cref{chap:3} to this statement, working through the components in order.

\paragraph{Entities.} The nouns in the problem that name classes of objects
are \emph{warehouses} (or \emph{locations}) and \emph{customers}. Each
becomes a set.

\begin{symboltable}{lL}
Symbol & Description \\ \midrule
\(F\) & Potential warehouse locations the company is considering. \\
\(C\) & Customers who must be served. \\
\end{symboltable}

\paragraph{Parameters.} The data known before any decision is made consists of two cost families.

\begin{symboltable}{lllL}
Symbol & Indices & Domain & Description \\ \midrule
\(f_i\) & \(i \in F\) & \(\mathbb{N}_0\) & Fixed cost of opening a warehouse at location \(i\). \\
\(c_{i,j}\) & \(i \in F,\; j \in C\) & \(\mathbb{N}_0\) & Cost of shipping from warehouse \(i\) to customer \(j\) \emph{if} that assignment is made. \\
\end{symboltable}

\paragraph{Variables.} There are two atomic decisions to make, whether
to open each warehouse, and whether each warehouse serves each
customer, giving one variable family per decision.

\begin{symboltable}{lllL}
Symbol & Indices & Domain & Description \\ \midrule
\(y_i\) & \(i \in F\) & \(\{0, 1\}\) & \lstinline!1! iff the warehouse at location \(i\) is opened. \\
\(x_{i,j}\) & \(i \in F,\; j \in C\) & \(\{0, 1\}\) & \lstinline!1! iff customer \(j\) is served by warehouse \(i\). \\
\end{symboltable}


\paragraph{Constraints.} The model has two families.

\emph{Assignment.} Every customer is served by exactly one warehouse.

\[
\sum_{i \in F} x_{i,j} \;=\; 1 \qquad \forall j \in C.
\]

\emph{Linking.} A warehouse may serve a customer only if the warehouse
is open.

\[
x_{i,j} \;\le\; y_i \qquad \forall i \in F, \; j \in C.
\]

If \(y_i = 0\) (warehouse closed), the inequality forces \(x_{i,j} = 0\)
for every customer \(j\). That warehouse cannot serve anyone. If
\(y_i = 1\) (warehouse open), the inequality is \(x_{i,j} \le 1\), which
the variable's domain already implies, so the constraint is silent. One
line captures the implication ``closed → cannot serve'' with no
implication arrow on paper.

\paragraph{Objective.} Minimize total cost: the fixed cost of every
warehouse opened, plus the shipping cost of every customer assignment
made.

\[
\min \quad \sum_{i \in F} f_i \, y_i \;+\; \sum_{i \in F} \sum_{j \in C} c_{i,j} \, x_{i,j}.
\]

That is the full mathematical model: two sets, two parameter families,
two variable families, two constraint families, one objective. The
CP-SAT translation is line-for-line:

\begin{lstlisting}[language=Python]
from ortools.sat.python import cp_model

# 1. Data
F = [1, 2, 3]                                  # potential locations
C = [1, 2, 3]                                  # customers
f = {1: 100, 2: 200, 3: 150}                   # opening cost f[i]
c = {1: {1: 10, 2: 20, 3: 30},                 # shipping cost c[i][j]
     2: {1: 15, 2: 25, 3: 35},
     3: {1: 20, 2: 30, 3: 40}}

model = cp_model.CpModel()

# 2. Variables
y = {i: model.new_bool_var(f"y_{i}") for i in F}
x = {(i, j): model.new_bool_var(f"x_{i}_{j}") for i in F for j in C}

# 3. Constraints
# Exactly one warehouse per customer:
for j in C:
    model.add(sum(x[i, j] for i in F) == 1)
# Only open warehouses can serve:
for i in F:
    for j in C:
        model.add(x[i, j] <= y[i])

# 4. Objective
model.minimize(
    sum(f[i] * y[i] for i in F) +
    sum(c[i][j] * x[i, j] for i in F for j in C)
)

# 5. Solve
solver = cp_model.CpSolver()
status = solver.solve(model)

# 6. Read the answer
if status in (cp_model.OPTIMAL, cp_model.FEASIBLE):
    opened = [i for i in F if solver.value(y[i])]
    served_by = {
        j: next(i for i in F if solver.value(x[i, j])) for j in C
    }
    print("opened warehouses:", opened)
    print("customer assignment:", served_by)
    print("total cost:", solver.objective_value)
else:
    print("no feasible solution")
\end{lstlisting}

\paragraph{Variables.} The code declares two families, mirroring the two
variable rows in the math model. \(y\) is one Boolean per potential location (\emph{is
this warehouse opened?}) created by
\lstinline!model.new_bool_var(name)! and corresponding
to \(y_i \in \{0, 1\}\) on paper. The \lstinline!name! is
metadata for debugging output and does not affect the search. \(x\) is
one Boolean per (location, customer) pair (\emph{does this warehouse
serve this customer?}) keyed by a tuple, corresponding to
\(x_{i,j} \in \{0, 1\}\). The dictionary keying is a Python convenience
that lets us write \lstinline!x[i, j]! later instead of
indexing into a flattened list.

Each paper variable here lifts directly into a CP-SAT variable with the
same domain and indexing. In larger models the correspondence is looser
(some paper variables are dropped as redundant, others are split into
auxiliary solver variables) but for this example the mapping is
direct.

\paragraph{Constraints.} The two families translate one-to-one into
Python. The assignment constraint, \emph{each customer gets exactly one
warehouse}, is one linear equality per customer:

\begin{lstlisting}[language=Python]
model.add(sum(x[i, j] for i in F) == 1)
\end{lstlisting}

Python's built-in \lstinline!sum(...)! builds a CP-SAT
linear expression by overloading addition on variables, and
\lstinline!model.add(... == 1)! registers the
linear-equality constraint. The \lstinline!==! is
\emph{building a constraint}, not evaluating to a Python bool. This
operator-overloading trick is at the heart of how CP-SAT's Python API
reads almost like the math: every arithmetic and comparison operator on
variables yields a CP-SAT object that the model collects.

The linking constraint, \emph{only open warehouses can serve}, is one
inequality per (location, customer) pair:

\begin{lstlisting}[language=Python]
model.add(x[i, j] <= y[i])
\end{lstlisting}

This is the line whose semantics we noted on paper: when \(y_i = 0\) the
right-hand side forces \(x_{i,j} = 0\); when \(y_i = 1\) the inequality
is slack. The Python expression is exactly the math: CP-SAT accepts
the linear inequality directly, with no extra machinery.

\paragraph{Objective.} \lstinline!model.minimize(expr)!
registers \lstinline!expr! as the objective and tells the
solver to minimize it. The expression is the same arithmetic object the
math model wrote: opening costs summed over locations, plus shipping
costs summed over (location, customer) pairs. To maximize instead, the
corresponding method is \lstinline!model.maximize!.

\paragraph{Solve.} \lstinline!solver = cp_model.CpSolver()!
instantiates the wrapper, and
\lstinline!status = solver.solve(model)! runs the search
and returns a status code. For a first model the two outcomes to know
are \lstinline!OPTIMAL! (a solution was found and proven optimal) and
\lstinline!FEASIBLE! (a solution was found, but optimality was not
proven, on hard problems this is a first-class outcome, often the
answer you ship). The remaining codes (\lstinline!INFEASIBLE!,
\lstinline!UNKNOWN!, \lstinline!MODEL_INVALID!) and time-limit handling
are covered in \Cref{chap:8}.

\paragraph{Read the answer.} Three pieces of information are worth
reading from the solver after a successful run.

\lstinline!solver.value(var)! returns the assigned value
of each variable as an integer: \lstinline!0! or
\lstinline!1! for Booleans, the assigned integer for an
integer variable. The companion
\lstinline!solver.boolean_value(literal)! returns the
same information as a Python \lstinline!bool! and is the
cleaner choice when downstream code wants
\lstinline!True!/\lstinline!False! directly.
Reading the answer through the \emph{solver} (not the \emph{variable})
is important either way: the variable object is just a handle; the
value lives in the solver's internal state after the search.

\lstinline!solver.objective_value! returns the objective
value of the solution: the optimum when status is
\lstinline!OPTIMAL!, the best-found value when status is
\lstinline!FEASIBLE! (where, again, optimality has not
been proven). For models without an objective there is nothing to read
here.

The two derived quantities, \lstinline!opened! and
\lstinline!served_by!, are simple comprehensions over
\lstinline!solver.value(...)!. They are not part of the
solver's output per se; they are how the program turns the variable
values into a domain-readable form.

Run the program, and you get something like:

\begin{lstlisting}
opened warehouses: [1]
customer assignment: {1: 1, 2: 1, 3: 1}
total cost: 160
\end{lstlisting}

In this trivial instance, opening only warehouse 1 (the cheapest at
100 to open, with shipping costs 10/20/30 to the three customers)
costs \lstinline!100 + 60 = 160!, which beats every other
configuration. CP-SAT proves this is optimal in essentially zero time.

In a real workflow, the solver's output is run through the verifier from
\Cref{chap:4} before being trusted: build a
\lstinline!Solution! object from the
\lstinline!solver.value(...)! calls, then call
\lstinline!evaluate(instance, sol)!. If the verifier
accepts the result, the encoding has produced a feasible solution and
the costs match. If it rejects, a bug in the encoding has just been
caught. The model here is short enough to read end to end, so we skip
the step.

\section{When the Natural Formulation Does Not Translate}\label{sec:5.5}

The facility-location model translated line-for-line from paper to code.
Not every problem is that smooth. Consider the \textbf{traveling
salesman problem} (TSP): given \(n\) cities and pairwise distances
\lstinline!cost[i][j]!, find a shortest tour visiting
every city exactly once and returning to the start.

The natural way to think about a tour is as a permutation: city
\(x_i\) is visited at position \(i\). And indeed, CP-SAT provides an
\lstinline!AllDifferent! constraint that handles exactly
this:

\begin{lstlisting}[language=Python]
x = [model.new_int_var(0, n - 1, f"x_{i}") for i in range(n)]
model.add_all_different(x)
\end{lstlisting}

So far so good. But the objective is
\lstinline!cost[x_0][x_1] + cost[x_1][x_2] + ...!,
and \(x_0\), \(x_1\) are \emph{variables}, not constants. We cannot
index into a Python list with a CP-SAT variable. Browsing the constraint
catalogue, we find \lstinline!add_element!, which
performs variable-indexed lookups, but only into one-dimensional
arrays. For a two-dimensional distance matrix, we would need nested
element constraints, each requiring auxiliary variables. The natural
formulation has no clean translation into CP-SAT.

But browsing further, we find \lstinline!add_circuit!,
a constraint that enforces a Hamiltonian cycle over a set of arcs. This
solves the problem directly, if we remodel. Instead of ``which city at
each position,'' ask ``which arcs are used?'' A Boolean per city pair
says whether the tour travels directly from \(i\) to \(j\):

\begin{lstlisting}[language=Python]
x = {(i, j): model.new_bool_var(f"x_{i}_{j}")
     for i in range(n) for j in range(n) if i != j}

model.add_circuit([(i, j, x[i, j]) for i, j in x])
model.minimize(sum(cost[i][j] * x[i, j] for i, j in x))
\end{lstlisting}

The objective is now a plain linear sum:
\lstinline!cost[i][j]! is a constant per arc, no variable
indexing needed. And \lstinline!add_circuit! handles
subtour elimination internally.

The natural description of a problem and the representation that fits
CP-SAT can be quite different things. When a formulation does not
translate cleanly, the answer is often not to force it through but to
look at what the solver \emph{can} express and remodel accordingly. The
constraint catalogue in \Cref{chap:7} is that vocabulary; it is worth
knowing what is in it before you start modeling a new problem.

\section{The Shape of a CP-SAT Program}\label{sec:5.6}

Every CP-SAT program follows the same six-step shape, with the middle
four steps directly mirroring the recipe from \Cref{chap:3}:

\begin{enumerate}
\item
  \textbf{Build the model object.}
  \lstinline!model = cp_model.CpModel()!.
\item
  \textbf{Add variables.} Calls to
  \lstinline!model.new_bool_var(...)!,
  \lstinline!model.new_int_var(...)!, or the other
  variable factories, typically one factory call per variable family
  in the paper model, plus any auxiliary variables the encoding needs
  (Step 3 of the recipe).
\item
  \textbf{Add constraints.} Calls to
  \lstinline!model.add(...)! for linear and logical
  constraints, and direct method calls
  (\lstinline!model.add_all_different(...)!,
  \lstinline!model.add_circuit(...)!) for global
  constraints. A simple paper constraint family often becomes a single
  call; a more complex one may need several
  \lstinline!add(...)! calls and helper variables to
  encode (Step 4 of the recipe).
\item
  \textbf{Add an objective} (if any).
  \lstinline!model.minimize(expr)! or
  \lstinline!model.maximize(expr)!. The expression is the
  same one the math model wrote (Step 5 of the recipe). Pure feasibility
  models (\emph{is there any valid configuration?}) skip this
  step.
\item
  \textbf{Solve.} \lstinline!solver.solve(model)! returns
  a status code.
\item
  \textbf{Read the answer.} \lstinline!solver.value(var)!
  for each variable of interest,
  \lstinline!solver.objective_value! for the objective,
  conditional on a successful status.
\end{enumerate}

The facility-location example used the simplest version of every piece;
the TSP example showed that not every problem is that smooth and that
knowing CP-SAT's vocabulary matters. The chapters that follow provide
that vocabulary. \Cref{chap:6} catalogs the variable types: integers with
bounded domains, intervals for scheduling, and several others. \Cref{chap:7}
catalogs the constraint families, including the global constraints
(\lstinline!add_all_different!,
\lstinline!add_circuit!,
\lstinline!add_no_overlap!,
\lstinline!add_cumulative!) that make CP-SAT what it is.
\Cref{chap:8} covers objectives, status codes, and time-limit handling.
\Cref{chap:9} puts the whole workflow loop together.

For now, you have the shape. The rest of the manuscript fills it in.
